% !TeX root = main.tex
\documentclass[10pt, a4paper]{article} % 1. 改为标准的 article，使用 a4paper
\usepackage{multicol}
\usepackage{geometry} % 虽然你用了手动设置，但引入这个包通常更稳
\usepackage{times}    % 2. 英文论文通常使用 Times New Roman 字体
\usepackage{enumitem}
\usepackage{indentfirst}

%%%%% 页面边距设置 (保留了你的原始数值) %%%%%%
\setlength{\textheight}{\paperheight}
\setlength{\topmargin}{-15.4mm}       % 3. 将 truemm 改为标准的 mm
\addtolength{\topmargin}{-\headheight}
\addtolength{\topmargin}{-\headsep}
\addtolength{\textheight}{-20mm}
\setlength{\textwidth}{\paperwidth}
\setlength{\oddsidemargin}{-10.4mm}
\setlength{\evensidemargin}{\oddsidemargin}
\addtolength{\textwidth}{-30mm}

%%%%% 行间距设置 %%%%%
\renewcommand{\baselinestretch}{.95}

\pagestyle{empty}

%%%%% 标题部分 %%%%%
\title{\textbf{Extend L* algorithm to Learn Syntactic Monoid}} % 加粗标题
\author{Haotian Zhang (Yuen/Nakazawa lab)}     % 英文作者名
\date{}

\begin{document}
\maketitle
\thispagestyle{empty}

%%%%% 正文分栏开始 %%%%%
\begin{multicols}{2} % 注意：通常学术双栏是 {2}，如果你确认要单栏请改成 {1}

\section{Introduction}
\subsection{L* Algorithm}
    The $L^*$ algorithm, introduced by Dana Angluin in 1987, is a seminal algorithm for learning regular languages via queries. \cite{angluin87}
    It operates based on an interaction between two entities: a Learner and a Teacher. The Learner's objective is to construct the minimal Deterministic Finite Automaton (DFA) for an unknown regular language $L$.
    The learning process relies on two specific types of queries:
\begin{enumerate}[leftmargin=*]
    \item \textbf{Membership Queries:} The Learner asks whether a specific string $w$ belongs to $L$. The Teacher responds with ``yes'' or ``no''.
    \item \textbf{Equivalence Queries:} The Learner proposes a hypothesis automaton $A$. The Teacher either confirms that $A$ is correct or provides a counterexample---a string that is accepted by $L$ but rejected by $A$ (or vice versa).
\end{enumerate}

    The Learner organizes these responses into an \textbf{observation table} $(I,S,T)$, which $I$ is the candidate state set, $S$ is the test suffix set, $T$ is a bool function representing the mapping in table. By maintaining the Closedness and separablity and properly dealing with counterexamples in $O(log|w|)$($|w|$ is the length of counterexample) \cite{Worrell:DFALearning} to iteratively refine this table, the Learner updates its hypothesis until it converges on the correct minimal DFA. Actually this is based on the Myhill-Nerode congruence: 
\begin{equation}
    u \equiv_L v \iff \forall w \in \Sigma^*, uw \in L \Leftrightarrow vw \in L
\end{equation}

\subsection{Syntactic Monoid}
Syntactic Monoid is the monoid with the minimum order that recognizes a regular language $L$. The underlying set is the quotient set by the syntactic congruence of $L$:
\begin{equation}
    u \sim_L v \iff \forall x,y \in \Sigma^*, xuy \in L \Leftrightarrow xvy \in L
\end{equation}

And the operation is the concatenation of representatives. Syntactic monoid is considered to characterize certain properties of regular languages, such as star-freeness and aperiodicity. 
\subsection{My research motivation and goal}
As for now, we know how to learn the minimal DFA of a regular language using the L* algorithm in polynomial time related to the size of the minimal DFA and the length of the longest counterexample. 
However, if we want to transform the minimal DFA to the syntactic monoid, it is $O(n^n)$ time complexity, where $n$ is the number of states of the minimal DFA. \cite{holzer2004deterministic}
Because when we want to decide some properties of regular languages, such as star-freeness and aperiodicity, we have to compute the syntactic monoid first. So it is natural to ask whether we can learn the syntactic monoid directly. 

My research goal is to extend the L* algorithm to directly learn the syntactic monoid of a regular language. By adapting the learning aspects of the L* algorithm, I aim to improve L* Algorithm that can efficiently construct the syntactic monoid based on membership and equivalence queries. 

\section{Idea and contribution}
The observation table in the original L* algorithm is two-dimensional, focused on test suffix because of the definition of the Myhill-Nerode congruence. 
To adapt this structure for learning syntactic monoids, the plain idea is to extend it to a three-dimensional observation table. 
In a new structure $(I,P,S,T)$, adding $P$ as the test prefix set, and modifying the bool function $T$ to accommodate the syntactic congruence definition. 

And because of the definition of syntactic congruence, we need to adjust the Closedness property of the observation table. We need two-side Closedness and we want to modify the acceptor into a new model that can represent monoids, such as $M = (Q,\Sigma,l/r,\delta,q_0,F)$, where $l/r$ is the left/right concatenation. 
Then we can use the "equivalence under the test words" to generalize the syntactic congruence by induction. 

For counterexamples, we now use the plain idea that adding all the prefixes and suffixes of the counterexample into $P$ and $S$. Although this is not efficient, if the algorithm can work correctly, we can further optimize it later.

\section{Future work}
\begin{enumerate}[leftmargin=*]
    \item \textbf{Proof of correctness:} Using the improvement we introduced, we need to prove that the new algorithm can always get the correct syntactic monoid of the target regular language.
    \item \textbf{Computational cost and improvement:} Theoretically evaluate if the extension algorithm is a polynomial-time algorithm. Try to reduce useless information during the process to improve computational cost, for example, optimize the handling of counterexamples.
    \item \textbf{Implementation and Experiment:} Implement the extended algorithm and conduct experiments for problems such as the decidability of star-free languages.
\end{enumerate}

%%%%% 参考文献 %%%%%
\bibliographystyle{plain} 
\bibliography{references}

\end{multicols}
\end{document}